\documentclass{article}
\usepackage{amsmath, amssymb, amsthm}
\usepackage{hyperref}
\usepackage{geometry}
\geometry{margin=1in}

\title{Erd\H{o}s Problem \#1054}
\date{Last updated: 2025-09-28}
\author{Source: erdosproblems.com}

\begin{document}
\maketitle

\noindent\textbf{Status:} OPEN \quad \textbf{No prize}

\noindent\textbf{Tags:} number theory, divisors

\medskip
\noindent\textbf{Problem Statement:}

Let $f(n)$ be the minimal integer $m$ such that $n$ is the sum of the $k$ smallest divisors of $m$ for some $k\geq 1$.

Is it true that $f(n)=o(n)$? Or is this true only for almost all $n$, and $\limsup f(n)/n=\infty$?

\bigskip
\noindent\textbf{Remarks:}

A question of Erd\H{o}s reported in problem B2 of Guy's collection \cite{Gu04}. The function $f(n)$ is undefined for $n=2$ and $n=5$, but is likely well-defined for all $n\geq 6$ (which would follow from a strong form of Goldbach's conjecture).

The sequence of values of $f(n)$ is given by A167485 in the OEIS.

See also [468].

The strong claim that $f(n)=o(n)$ was disproved by Tao in the comments to [468], in which he proves that the upper density of $\{ n : f(n)\leq \delta n\}$ is $\ll \delta^2$.

\bigskip
\noindent\textbf{References:}

[Gu04] Guy, Richard K., Unsolved problems in number theory. (2004), xviii+437.

\bigskip
\noindent\small{Source: \url{https://www.erdosproblems.com/1054}}
\end{document}
